\begin{lemma}
\label{lemma-irreducible}
Let $R$ be a ring.
\begin{enumerate}
\item For a prime $\mathfrak p \subset R$ the closure
of $\{\mathfrak p\}$ in the Zariski topology is $V(\mathfrak p)$.
In a formula $\overline{\{\mathfrak p\}} = V(\mathfrak p)$.
\item The irreducible closed subsets of $\Spec(R)$ are
exactly the subsets $V(\mathfrak p)$, with $\mathfrak p \subset R$
a prime.
\item The irreducible components (see Topology,
Definition \ref{topology-definition-irreducible-components})
of $\Spec(R)$ are  exactly the subsets $V(\mathfrak p)$,
with $\mathfrak p \subset R$ a minimal prime.
\end{enumerate}
\end{lemma}

\begin{proof}
Note that if $ \mathfrak p \in V(I)$, then
$I \subset \mathfrak p$. Hence,
clearly $\overline{\{\mathfrak p\}} = V(\mathfrak p)$.
In particular $V(\mathfrak p)$ is the closure of
a singleton and hence irreducible.
The second assertion implies the third.
To show the second, let
$V(I) \subset \Spec(R)$ with $I$ a radical ideal.
If $I$ is not prime, then choose $a, b\in R$, $a, b\not \in I$
with $ab\in I$. In this case $V(I, a) \cup V(I, b) = V(I)$,
but neither $V(I, b) = V(I)$ nor $V(I, a) = V(I)$, by
Lemma \ref{lemma-Zariski-topology}. Hence $V(I)$ is not
irreducible.
\end{proof}